\section{Frequency-based techniques}

\subsection{EM coil current control}
\begin{figure}[H]
    \centering
    \resizebox{0.95\textwidth}{!}{%
    \begin{tikzpicture}[
        auto,
        node distance=1.2cm and 1.2cm,
        >=Latex,
        block/.style={
            draw,
            rectangle,
            minimum width=1.5cm,
            minimum height=1.0cm,
            align=center
        },
        sum/.style={
            draw,
            circle,
            minimum size=8mm
        },
        input/.style={coordinate},
        output/.style={coordinate},
        satblock/.style={
            draw,
            rectangle,
            minimum width=2.1cm,
            minimum height=1.45cm,
            inner sep=0pt,
            %label={[font=\scriptsize]below:$\mathrm{sat}(\cdot)$},
            path picture={
                % assi centrati in 0
                \draw[->]
                    ([xshift=-0.75cm,yshift=0cm]path picture bounding box.center) --
                    ([xshift= 0.78cm,yshift=0cm]path picture bounding box.center);
                \draw[->]
                    ([xshift=0cm,yshift=-0.48cm]path picture bounding box.center) --
                    ([xshift=0cm,yshift= 0.50cm]path picture bounding box.center);
                % caratteristica di saturazione centrata in 0
                \draw[thick]
                    ([xshift=-0.65cm,yshift=-0.32cm]path picture bounding box.center) --
                    ([xshift=-0.25cm,yshift=-0.32cm]path picture bounding box.center) --
                    ([xshift= 0.25cm,yshift= 0.32cm]path picture bounding box.center) --
                    ([xshift= 0.65cm,yshift= 0.32cm]path picture bounding box.center);
                % livelli di saturazione
                \draw[dashed]
                    ([xshift=-0.68cm,yshift= 0.32cm]path picture bounding box.center) --
                    ([xshift= 0.68cm,yshift= 0.32cm]path picture bounding box.center);
                \draw[dashed]
                    ([xshift=-0.68cm,yshift=-0.32cm]path picture bounding box.center) --
                    ([xshift= 0.68cm,yshift=-0.32cm]path picture bounding box.center);
                % etichette
                \node[font=\tiny, anchor=west] at
                    ([xshift=0.70cm,yshift= 0.32cm]path picture bounding box.center)
                    {$+V_{\max}$};
                \node[font=\tiny, anchor=west] at
                    ([xshift=0.70cm,yshift=-0.32cm]path picture bounding box.center)
                    {$-V_{\max}$};
                \node[font=\tiny, anchor=north east] at
                    ([xshift=0cm,yshift=0cm]path picture bounding box.center)
                    {$0$};
            }
        },
        font=\small
    ]
    % Nodi
    \node[input] (input) {};
    \node[sum, right=of input] (sum) {$+$};
    \node[block, right=of sum] (controller) {$PI(s)$};
    \node[block, right=of controller] (amp) {$K_g$};
    \node[satblock, right=of amp] (sat) {};
    \node[block, right=of sat] (plant) {$\displaystyle \frac{1}{Ls+R}$};
    % output più lontano -> freccia più lunga
    \node[output, right=2.0cm of plant] (output) {};
    % Connessioni
    \draw[->] (input) -- node[pos=0.2, above] {$i_{\mathrm{ref}}(t)$} (sum);
    \draw[->] (sum) -- node[above] {$e_i(t)$} (controller);
    \draw[->] (controller) -- node[above] {$u_c(t)$} (amp);
    \draw[->] (amp) -- node[above] {$v_c(t)$} (sat);
    \draw[->] (sat) -- node[above] {$v_a(t)$} (plant);
    \draw[->] (plant) -- node[pos=0.6, above] {$i(t)$} (output);
    % Retroazione
    \draw[->] (output) |- ++(0,-1.9) -| node[pos=0.95, left] {$-$} (sum);
    \end{tikzpicture}%
    }
    \caption{Current control loop}
    \label{fig:block_current_control_sat}
\end{figure}


The first control task is the design of an inner current loop. The open-loop transfer function from the amplifier command voltage to the coil current is obtained from the RL model:
\begin{equation}
    G_i(s) = \frac{I(s)}{U(s)} = \frac{K_g}{Ls + R}.
\end{equation}

A PI controller is selected to obtain zero steady-state error for constant current references:
\begin{equation}
    R_i(s) = K_{p,i} + \frac{K_{i,i}}{s}.
\end{equation}

The resulting loop transfer function is
\begin{equation}
    L_i(s) = R_i(s)G_i(s).
\end{equation}

The closed-loop current transfer function is
\begin{equation}
    H_i(s) = \frac{R_i(s)G_i(s)}{1+R_i(s)G_i(s)}.
\end{equation}

\subsection{Position control}

After closing the current loop, the outer plant can be approximated by the transfer function from current reference variation $\Delta i_{ref}$ to position variation $\Delta x$. The controller is designed using Bode diagrams, root locus or direct loop shaping.

The design objective is to stabilize the unstable mechanical plant while keeping sufficient robustness margins. The resulting open-loop transfer function is
\begin{equation}
    L_x(s) = R_x(s)G_x(s),
\end{equation}
where $R_x(s)$ is the position controller and $G_x(s)$ is the linearized outer-loop plant.

\subsection{Controller design}

The selected controller structure is reported here. Replace this placeholder with the final transfer function used in the experiments:
\begin{equation}
    R_x(s) = \frac{N_R(s)}{D_R(s)}.
\end{equation}

The design is validated by checking gain margin, phase margin and crossover frequency:
\begin{lstlisting}[style=matlabstyle,caption={Frequency-based validation commands.}]
Lx = Rx * Gx;
margin(Lx);
Tcl = feedback(Lx, 1);
step(Tcl);
\end{lstlisting}

\subsection{Simulation testing}

The controller is first tested in simulation on the linear model and then on the nonlinear model. The comparison allows us to evaluate the validity of the linear approximation and the effect of saturation.

% \begin{figure}[H]
%     \centering
%     \includegraphics[width=0.85\textwidth]{fb_step_simulation.pdf}
%     \caption{Simulation response of the frequency-based position controller.}
%     \label{fig:fb_step_simulation}
% \end{figure}

\subsection{Real system implementation and validation}

The final controller is implemented in Simulink and tested on the physical setup. The experimental response is compared with the simulated one.

\subsubsection{Time validation}

Report here the measured step response, overshoot, settling time, steady-state error and current effort.

\subsubsection{Frequency validation}

Report here the measured sinusoidal responses and the comparison with the predicted frequency response.
